\documentclass[11pt]{article}
\usepackage[margin=0.9in]{geometry}
\usepackage{amsmath,amssymb,amsthm,booktabs,mathrsfs}
\usepackage[T1]{fontenc}
\usepackage{lmodern,microtype}
\usepackage[colorlinks=true,linkcolor=blue,urlcolor=blue]{hyperref}
\newtheorem{theorem}{Theorem}[section]
\newtheorem{lemma}[theorem]{Lemma}
\newtheorem{proposition}[theorem]{Proposition}
\newtheorem{corollary}[theorem]{Corollary}
\theoremstyle{remark}
\newtheorem{remark}[theorem]{Remark}
\newcommand{\norm}[1]{\left\|#1\right\|}
\title{Toward a four-dimensional Yang--Mills construction\\
\large Notebook 28: a fixed analytic ball and high-degree stability}
\author{Research working document}
\date{17 October 2026}
\begin{document}
\maketitle
\begin{abstract}
The generic Taylor-at-zero transport estimate shrinks the background
domain at each use. For the composite macrostep of Notebook 27,
we avoid this loss by composing inherited activities around a
bounded local core and expanding only the distant tail. The
pre-gauge representative is small by an inverse macro ratio.
A local analytic map supplies strict joint background/fluctuation
families for the entire inherited insertion. Cauchy estimates in
the fluctuation scaling parameter preserve its small factor in
the anchored activity norm. Classical precision and base
interactions are substituted into the stationary background
before Gaussian integration, using their independent classical
domains and finite support reserves. The Gaussian output is then
already a function of the next background.
A fixed sequence of parameter choices yields one inherited
small-field update on the same positive input and output
analytic radius, uniformly in preceding refinement and volume.
The general activity bound is not contractive. A separate
estimate proves an inverse-ratio contraction for background
transport when every incoming coefficient has vanishing
field Taylor jets through degree four. A bounded Taylor
subtraction produces such a remainder. The same estimate,
combined with a two-input Gaussian logarithm comparison,
contracts the full quantum response to input differences
whose coefficientwise low jets agree. Its high-degree
output projection can have Lipschitz constant at most
one-half. The full update can still regenerate low jets,
and this subtraction does not supply a gauge-invariant
extraction of them. These results close the stated domain
and support interfaces of a finite step and supply a
conditional quantum contraction, while leaving uniform
quantum amplitudes, coupling flow and the remaining
sectors of the existence argument open.
\end{abstract}

All identities below are first finite-regulator identities on
common small domains. Enlarged-domain identities are identified
as analytic continuations of those germs. This does not by itself
identify a globally continued homogeneous sector with the full
Wilson measure or control histories containing large-field
regions. The primitive ratio, composite blocking map, gauge
choice, retained link/color cutoff and normalization are those
constructed in Notebook 27.

